\section{引言：从中子电流档案到可追溯的工程证据}

核电厂的仪表与控制系统持续产生用于运行监视、状态评估和维护决策的测量记录。在线监测的目标并不只是发现数值偏离，更是以可追溯的方式评估仪表通道、过程和部件的状态，并为后续核查提供依据。国际原子能机构的在线监测指南将仪表通道监测与过程/部件状态监测分别作为核电厂性能改进的重要组成部分 \citep{iaea2008olm1,iaea2008olm2}；相关综述也指出，在线监测需要将测量信息、诊断方法与维护决策联系起来。
\citep{hashemian2011online}。

在核反应堆的测量指标中，中子电流是在线监测中的重要测量对象：中子噪声分析可非侵入地研究在反应堆运行和中子测量系统有关的现象 \citep{pazsit2010noise}，数值模拟、建模和验证工作也服务于反应堆监测与安全研究 \citep{demaziere2022cortex}。本文聚焦核电厂运行环境中积累的无标签多通道中子电流数据记录。这些记录包含时间变化、通道间关系和信号完整性线索，却通常缺少统一的样本级故障真值。

现有中子探测器验证、噪声监测和核电厂故障诊断研究，通常在预设信号构造、任务边界或故障场景下开展，因此能够以明确的分类或识别目标评价方法 \citep{tambouratzis2020neutron,tambouratzis2020decomposition,elbordany2024fault}；检索增强生成则可利用外部资料组织异构技术知识 \citep{lewis2020rag}。两条路线都不能直接解决无标签中子电流记录的解释问题：前者缺少可直接归类的统一故障真值，后者不能替代对原始测量的分析；在安全相关工程场景中，来源、定位和适用范围不受控制的流畅解释也不能成为可靠依据。

若要将无标签数据中的局部波形可靠归为某一设备故障，仍需结合运行工况、维护历史和核工业专业知识；为全体长序列逐点补充这些上下文和专家标签，在数据规模上并不现实。本文因此研究如何在没有统一故障标签且不越过工程复核边界的条件下，将原始中子电流数据转化为工程师可审查的观察，并使后续技术解释可追溯至适当的权威资料。这个问题既不是具有已知类别、可报告分类准确率的监督学习基准，也不是单纯的异常检测或文档问答，而是要求在测量与证据之间建立明确、可复现的中间层。

本文据此提出\textbf{可追溯的测量—证据接口}（图~\ref{fig:introduction-overview}），以可复核观察为中间层，将确定性分析的测量侧与受来源策略约束的证据侧连接起来；工程师复核保留在接口之外。

\begin{figure*}[t]
  \centering
  \includegraphics[width=0.62\textwidth]{fig01_overview}
  \caption{测量—证据接口的研究总览（第二版占位图）。原始中子电流记录经确定性分析形成可回溯观察与结构化摘要；摘要在来源策略约束下连接至可定位的技术证据，并组织为供工程师复核的候选解释。正式版本将按第三版术语重绘。}
  \label{fig:introduction-overview}
\end{figure*}

首先，确定性分析工具从原始时序数据中提取数据健康、突刺、零值、滞后、共变和极值等\textbf{观测客观信息}，并将输入、参数和结构化结果保存为派生产物。稳定的观测字段直接构成紧凑告警字段，作为受控的检索描述，而不是判断性的故障标签；测量侧语言模型仅依据这些已保存的结果组织供人工阅读的描述报告。随后，证据侧以紧凑告警字段检索带有来源元数据的可定位段落，并由来源策略约束预设权威资料在前 10 个检索结果中所占的比例（下文简称“权威资料可见性”）；证据侧语言模型据此组织带引文的候选解释。两侧语言模型均不确认物理机理，也不替代工程师作判断。本文不把最高语义相似度视为充分标准。特别是在由核心权威资料和同领域补充资料共同构成的资料库中，语义上更接近的补充文档可能使前 10 个结果中预设权威资料的占比下降。因此，技术证据的获取并非由单一检索配置决定；查询表述、语言、编码器、分块方式和来源策略会共同影响检索结果中权威资料所占比例，因而需要作为证据来源约束问题加以考察。

本文首先提出面向无统一故障真值中子电流档案的可追溯测量—证据接口，将确定性脚本生成的观测与受来源策略约束的技术证据连接起来，为人工复核提供候选解释。其次，围绕该接口构建检索治理评价，系统考察资料范围、查询表述、语言、编码器和分块策略对技术证据获取及前 10 个结果中预设权威资料占比的影响。本文始终将工程判断保留在人类复核环节。
